\newpage
\section{ État de l'art}
\label{chap:etat_art}

Ce chapitre situe le projet CYBEL par rapport aux travaux et solutions
existants. Il présente les fondements théoriques de la robotique de service
ainsi que les solutions logicielles et matérielles comparables, afin de
mettre en évidence les limites des approches actuelles et de justifier les
choix architecturaux retenus.

\subsection{Fondements théoriques}
\label{subsec:fondements_theoriques}

La conception de la plateforme s'appuie sur plusieurs corpus de référence :

\begin{itemize}
    \item \textbf{Architecture logicielle robotique} : le modèle ROS
    (topics, services, actions) proposé par Quigley et al. constitue le
    paradigme de communication analysé dans ce projet.

    \item \textbf{Navigation autonome} : les concepts de cartes
    d'occupation et de planification (costmaps, DWA) décrits dans
    \textit{Probabilistic Robotics} (Thrun, Burgard \& Fox, 2005)
    servent de base à l'interprétation du module \texttt{move\_base}.

    \item \textbf{Protocole rosbridge} : le standard JSON de communication
    ROS (subscribe, publish, call\_service) est utilisé comme référence pour
    la reconstruction des échanges observés.

    \item \textbf{Interaction homme-robot (HRI)} : les principes de
    supervision en temps réel et de sécurité opérationnelle (priorité à
    l'arrêt d'urgence) guident la conception de l'interface CYBEL.

    \item \textbf{Sécurité IoT} : les recommandations de l'OWASP IoT Top 10
    permettent d'analyser les risques liés aux canaux non authentifiés
    observés sur le système.
\end{itemize}

\subsection{État des solutions existantes}
\label{subsec:solutions_existantes}

Les solutions suivantes représentent les principales approches disponibles
dans l'écosystème robotique et du robot étudié.

\begin{table}[h!]
\centering
\begin{tabular}{|p{0.28\textwidth}|p{0.38\textwidth}|p{0.24\textwidth}|}
\hline
\textbf{Solution} & \textbf{Description} & \textbf{Accessibilité} \\
\hline
Application propriétaire (upper body Android) &
Interface tactile constructeur pour accueil et navigation &
Fermée, sans API ni documentation \\
\hline
Interface web embarquée (\texttt{:8082}) &
Outil de gestion des cartes SLAM &
Accessible localement, non documenté \\
\hline
Interface de debug (\texttt{:8088}) &
Interface interne de diagnostic &
Non documentée, usage inconnu \\
\hline
rosbridge / RViz / Foxglove &
Outils ROS de visualisation et contrôle &
Open-source mais génériques \\
\hline
SDK robots commerciaux (Temi, Pepper, etc.) &
SDK propriétaires pour robots de service &
Documentés mais verrouillés \\
\hline
\end{tabular}
\caption{Panorama des solutions existantes comparables}
\label{tab:solutions_existantes}
\end{table}

\subsection{Analyse comparative}
\label{subsec:benchmark}

\begin{table}[h!]
\centering
\begin{tabular}{|p{0.28\textwidth}|p{0.20\textwidth}|p{0.20\textwidth}|p{0.20\textwidth}|}
\hline
\textbf{Critère} &
\textbf{App. propriétaire} &
\textbf{RViz / Foxglove} &
\textbf{CYBEL} \\
\hline
Documentation &
Aucune &
Open-source &
Documenté (rapport + code) \\
\hline
Accès distant web &
Non &
Partiel &
Oui \\
\hline
Personnalisation scénarios &
Non &
Non &
Oui \\
\hline
Mode simulation &
Non &
Partiel &
Oui (mock) \\
\hline
Navigation interactive &
Oui &
Oui &
Oui + logique métier \\
\hline
Interaction vocale &
Oui &
Non &
Partiel / évolutif \\
\hline
\end{tabular}
\caption{Benchmark comparatif des solutions}
\label{tab:benchmark_cybel}
\end{table}

\subsection{Limites des solutions existantes}
\label{subsec:limites}

Les solutions étudiées présentent plusieurs limites majeures :

\begin{itemize}
    \item Absence d'extensibilité des systèmes propriétaires
    \item Outils ROS de trop bas niveau pour des usages métier
    \item Aucune solution ne propose de mode déconnecté du robot
    \item Documentation constructeur insuffisante nécessitant rétro-ingénierie
\end{itemize}

\subsection{Synthèse}
\label{subsec:synthese_etat_art}

L'analyse met en évidence un manque de solutions adaptées à un usage
orienté \textbf{scénarios d'accueil personnalisés} et \textbf{supervision
haut niveau}. Ces limites justifient le développement de la plateforme
CYBEL, détaillée dans le chapitre suivant.

\subsection{Conclusion}
\label{subsec:conclusion_etat_art}

Ce chapitre a permis de situer le projet CYBEL par rapport aux fondements
théoriques de la robotique de service architecture ROS, navigation
autonome, protocole \texttt{rosbridge}, interaction homme-robot et sécurité
des objets connectés ainsi qu'aux solutions existantes, qu'elles soient
propriétaires (application Android constructeur) ou génériques (rosbridge,
RViz, Foxglove, SDK de robots commerciaux). L'analyse comparative a montré
qu'aucune de ces solutions ne répond pleinement aux besoins identifiés au
Chapitre 2 : extensibilité, logique métier orientée accueil, supervision
de haut niveau et mode de développement déconnecté du robot physique. Ces
limites confirment la pertinence du projet CYBEL et orientent directement
les choix de conception présentés au Chapitre 4, qui détaille l'analyse
fonctionnelle et l'architecture logicielle de la solution retenue.